\section*{الفصل \ref{ch:b3:environmental-toxicology} --- الكيمياء البيئية وعلم السموم}

\begin{solution}{exo:b3:environmental-toxicology:1}
للجزيء الثنائي الذرة المتجانس النواة اهتزاز واحد يُبقي عزم ثنائي القطب صفرًا،
والأرغون لا اهتزاز له: فلا يستطيع أي منها امتصاص الأشعة تحت الحمراء. ويمتص $\ce{CO2}$ عبر
نمط الانحناء فيه (وعبر تمططه اللامتناظر)، اللذين يُحدثان ثنائي قطب
متذبذبًا؛ ويقع الانحناء في وسط إصدار الأرض.
\end{solution}

\begin{solution}{exo:b3:environmental-toxicology:2}
$10^{0.1} = 1.26$: يرتفع $[\ce{H+}]$ بنسبة 26\,\%.
\end{solution}

\begin{solution}{exo:b3:environmental-toxicology:3}
$(2.0 + 0.50)\ \unit{mmol/L} \times \qty{100.1}{mg/mmol} = \qty{250}{mg/L}$ من $\ce{CaCO3}$.
\end{solution}

\begin{solution}{exo:b3:environmental-toxicology:4}
النيتروجين: $0.60/14.0 = \qty{0.043}{mmol/L}$؛ والفوسفور: $0.010/31.0 = \qty{3.2e-4}{mmol/L}$؛ ونسبة النيتروجين إلى الفوسفور $= 133$، أعلى بكثير من 16: فالفوسفور ينفد
أولًا ويحدّ من النمو.
\end{solution}

\begin{solution}{exo:b3:environmental-toxicology:5}
$L_0 = \mathrm{BOD_5}/(1 - \eu^{-0.23 \times 5}) = 170/0.683 = \qty{250}{mg/L}$.
\end{solution}

\begin{solution}{exo:b3:environmental-toxicology:6}
$\qty{4.40e-3}{L} \times \qty{0.0200}{mol/L} = \qty{8.80e-5}{mol}$ في \qty{0.1000}{L}: \omterm{def:b3:environmental-toxicology:alkalinity}{القلوية} $\qty{8.80e-4}{mol/L}$؛ وفي صورة $\ce{HCO3-}$ (\qty{61.0}{g/mol})،
\qty{53.7}{mg/L}.
\end{solution}

\begin{solution}{exo:b3:environmental-toxicology:7}
$K_d = 0.020 \times 300 = \qty{6.0}{L/kg}$. ونسبة الممتز إلى الكلي $= K_dm_s/(K_dm_s + V_w) = 9.0/(9.0 + 0.30) = 0.97$: أي 97\,\% ممتز.
\end{solution}

\begin{solution}{exo:b3:environmental-toxicology:8}
البنزين: $\log\mathrm{BCF} = 0.85 \times 2.13 - 0.70 = 1.11$، \omterm{def:b3:environmental-toxicology:bio}{وعامل التركيز الحيوي} نحو 13. ومركب PCB: $0.85 \times 6.67 - 0.70 = 4.97$، \omterm{def:b3:environmental-toxicology:bio}{وعامل التركيز الحيوي} نحو
\num{9e4}: فهو يتركز في الأسماك أكثر من البنزين بنحو سبعة آلاف مرة؛ وعند مثل هذه القيم للمعامل
$K_{ow}$ يقترب الانحدار من حد صلاحيته.
\end{solution}

\begin{solution}{exo:b3:environmental-toxicology:9}
$\qty{192}{mg/kg} \times \qty{70}{kg} = \qty{13}{g}$، أي نحو 150 فنجانًا. فالأنواع تختلف في الامتصاص والاستقلاب،
وقيمة $\mathrm{LD_{50}}$ ليست عتبة أمان؛ فالآثار الضارة تبدأ دونها بكثير.
\end{solution}

\begin{solution}{exo:b3:environmental-toxicology:10}
PNEC $= \qty{0.50}{mg/L}/1000 = \qty{0.50}{\micro g/L}$؛ و$\mathrm{RQ} = 0.20/0.50 = 0.40$، أقل من 1: لا ما يقلق عند هذا التعرض.
\end{solution}

\begin{solution}{exo:b3:environmental-toxicology:11}
العوالق/الماء: $\qty{0.015}{mg/kg}/\qty{2.0e-6}{mg/L} = \qty{7.5e3}{L/kg}$ (تركيز حيوي). الأسماك/العوالق: $0.30/0.015 = 20$؛
الطائر/الأسماك: $6.0/0.30 = 20$ (\omterm{def:b3:environmental-toxicology:bio}{تضخم حيوي} في كل خطوة)؛ الطائر/الماء: $\num{3e6}$.
\end{solution}

\begin{solution}{exo:b3:environmental-toxicology:12}
$\qty{1.0e9}{g}/\qty{64.1}{g/mol} = \qty{1.56e7}{mol}$ من $\ce{SO2}$، تعطي العدد نفسه من مولات $\ce{H2SO4}$: $\qty{1.53e9}{g}$، أي نحو \qty{1500}{t}. وهي
تحرر $\qty{3.12e7}{mol}$ من $\ce{H+}$؛ وعند pH يساوي 4.0، $[\ce{H+}] = \qty{1.0e-4}{mol/L}$: أي $\qty{3.1e11}{L}$ من المطر.
\end{solution}

\begin{solution}{pb:b3:environmental-toxicology:1}
\textbf{1.} ذوبانه في الأوكتانول أكبر من ذوبانه في الماء بعشرة آلاف مرة: فهو كاره للماء،
وسيمتز على المادة العضوية ويتراكم في الدهون.

\textbf{2.} $K_{oc} = 0.41 \times 10^4 = \qty{4100}{L/kg}$؛ $K_{sw} = f_{oc}K_{oc} = 0.040 \times 4100 = \qty{164}{L/kg}$.

\textbf{3.} $K_{fw} = 0.048 \times 10^4 = \qty{480}{L/kg}$.

\textbf{4.} لا: عند التوازن يحوي الهواء $10^{-4}$ من التركيز في الماء.

\textbf{5.} الجزيء المتعادل الكاره للماء يذوب في المادة العضوية، كما في
الأوكتانول؛ أما السطوح المعدنية فقطبية ومغطاة بالماء.

\textbf{6.} الرواسب، التي سعتها، $K_{sw}m_s$، هي الأكبر.

\textbf{7.} $V_w = \qty{1.0e9}{L}$؛ $K_{aw}V_a = \qty{1.0e8}{L}$؛ $K_{sw}m_s = \qty{1.64e10}{L}$؛ $K_{fw}m_f = \qty{4.8e6}{L}$.

\textbf{8.} $C_w = \qty{1.0e8}{mg}/\qty{1.75e10}{L} = \qty{5.7e-3}{mg/L}$، أي \qty{5.7}{\micro g/L}.

\textbf{9.} الماء \qty{5.7}{kg}، والهواء \qty{0.57}{kg}، والرواسب \qty{94}{kg}، والأسماك \qty{0.027}{kg}.

\textbf{10.} الهواء $\qty{5.7e-7}{mg/L}$؛ والرواسب \qty{0.94}{mg/kg}؛ والأسماك \qty{2.7}{mg/kg}.

\textbf{11.} التوازن بين كل الأحياز في آن واحد، دون تحلل، ودون
تدفق داخل أو خارج، مع أحياز جيدة المزج.

\textbf{12.} مع التحلل والتدفقات في حالة مستقرة (المستوى الثاني) تحدد المخزونَ
سرعاتُ الإدخال والفقد؛ ومع سرعات انتقال بين الأحياز (المستوى الثالث)
لا تبقى التراكيز في نسب التوازن، ويُثرى الحيز الذي
يتلقى الإدخال.

\textbf{13.} $k = \ln 2/\qty{60}{d} = \qty{0.0116}{d^{-1}}$.

\textbf{14.} $\eu^{-0.0116 \times 365} = 0.0145$: يبقى 1.5\,\%.

\textbf{15.} $\qty{5.7}{\micro g/L} \times 0.0145 = \qty{0.083}{\micro g/L}$.

\textbf{16.} التحلل غالبًا أبطأ في الرواسب الباردة العديمة الأكسجين، والمبيد الممتز
يتحرر عائدًا إلى الماء ببطء، فيُبقيه حاضرًا بعد أن يكون
الماء قد صفا منه.

\textbf{17.} روابط تقاوم الحلمأة والأكسدة والهجوم الميكروبي (الرابطة C--Cl
العطرية مثلًا)، وذوبانية ضعيفة في الماء، \omterm{def:b3:surfaces-catalysis:adsorption}{وامتزاز} يحميه من
الكائنات المحللة.

\textbf{18.} PNEC $= \qty{50}{\micro g/L}/10 = \qty{5.0}{\micro g/L}$.

\textbf{19.} $\mathrm{RQ} = 5.7/5.0 = 1.1$.

\textbf{20.} أعلى من 1 بقليل: هناك خطر مُشار إليه، ويجب تحسين التقييم (بيانات
سمية أفضل، وتراكيز مقيسة) أو خفض التعرض.

\textbf{21.} \qty{2.7}{mg/kg}، أي نحو ثلاثة أضعاف العتبة \qty{1.0}{mg/kg}: فلا ينبغي أكل الأسماك
حتى ينخفض التركيز.

\textbf{22.} التراكيز في الماء والرواسب والأسماك على مر الزمن وفي عدة نقاط،
للتحقق من النموذج ومتابعة الانخفاض.

\textbf{23.} استعمال كمية أقل من المبيد، ورشه بعيدًا عن المطر وعن الشاطئ،
وشرائط عازلة تحتجز الجريان السطحي، أو بديل أقل ثباتًا.

\textbf{24.} يقسم المستوى المقيس عديم الأثر ليغطي الفجوة بين بضعة أنواع
مخبرية ونظام بيئي كامل؛ ومعامل قدره 10 بدلًا من 1000 يغيّر
PNEC، والحكم، مئة ضعف.

\textbf{25.} \omterm{def:b3:environmental-toxicology:ecotoxicology}{حاصل الخطر} PEC/PNEC نحو 1.1، أي أعلى من 1 بقليل.
\end{solution}
